\subsection{Stabilization of the drift network $K$}

In the baseline specification, the drift network $K$ is modeled as a linear map
\begin{equation}
\mu(z)=Mz+N,
\label{eq:k_baseline}
\end{equation}
where $M \in \mathbb{R}^{d \times d}$ and $N \in \mathbb{R}^d$ are trainable parameters.
This formulation is simple and expressive, but it does not guarantee stable latent
dynamics under forward simulation. In particular, if the learned matrix $M$ has one or
more eigenvalues with non-negative real part, the deterministic part of the state
dynamics may fail to be mean-reverting, which can lead to unstable or explosive sample
paths.

To address this issue, I replace the unconstrained matrix $M$ by the structured
parameterization
\begin{equation}
M = -\left(V^\top V + \varepsilon I_d\right),
\label{eq:k_stable_matrix}
\end{equation}
where $V \in \mathbb{R}^{d \times d}$ is a trainable matrix, $I_d$ is the identity matrix
in dimension $d$, and $\varepsilon > 0$ is a small constant. The drift network is then
given by
\begin{equation}
\mu(z)= -\left(V^\top V + \varepsilon I_d\right)z + N.
\label{eq:k_stable_drift}
\end{equation}

This reparameterization guarantees that the drift matrix is negative definite. Since
$V^\top V$ is positive semidefinite, the matrix $V^\top V + \varepsilon I_d$ is strictly
positive definite for every $\varepsilon > 0$. Hence, for any nonzero vector
$x \in \mathbb{R}^d$,
\begin{equation}
x^\top M x
=
- x^\top \left(V^\top V + \varepsilon I_d\right)x
=
-\|Vx\|^2 - \varepsilon \|x\|^2
< 0.
\end{equation}
It follows that all eigenvalues of $M$ are real and strictly negative. The linear drift
therefore has an inherent pull toward an equilibrium level and cannot generate locally
explosive deterministic behavior.

The equilibrium point is obtained by solving $\mu(z^\ast)=0$, which yields
\begin{equation}
z^\ast = \left(V^\top V + \varepsilon I_d\right)^{-1}N.
\end{equation}
Equivalently, the drift can be written as
\begin{equation}
\mu(z)=-Q(z-z^\ast),
\qquad
Q:=V^\top V+\varepsilon I_d,
\end{equation}
which makes the mean-reverting structure explicit: deviations from the equilibrium state
$z^\ast$ are pulled back at a rate governed by the positive definite matrix $Q$.

The motivation for this modification was numerical instability observed when simulating
from the original unconstrained specification in \eqref{eq:k_baseline}. In that case,
stability depends entirely on the outcome of the optimization, since the learned matrix
$M$ may contain positive or weakly damped eigenvalues. By contrast, the parameterization
in \eqref{eq:k_stable_matrix} guarantees a stable linear drift for every value of the
trainable matrix $V$. This is particularly desirable in the present setting, where the
latent factors enter a continuous-time term-structure model and must remain well behaved
under forward simulation.

An additional advantage is that the stabilized parameterization does not increase the
number of trainable parameters. The original formulation uses $d^2$ parameters for $M$
and $d$ parameters for $N$, while the stabilized version uses $d^2$ parameters for $V$
and $d$ parameters for $N$. Thus, both specifications contain $d^2+d$ trainable
parameters. The extra matrix multiplication required to form $V^\top V$ is of the same
asymptotic order as the original linear mapping and is negligible in the present low-
dimensional setting.

In the implementation, $V$ is initialized orthogonally and $N$ is initialized at zero.
Orthogonal initialization provides a well-scaled starting point for optimization, while
the constant $\varepsilon$ ensures strict negative definiteness even in directions where
$V^\top V$ is close to singular. The resulting network preserves the linear form of the
original drift specification, but enforces stable and mean-reverting latent dynamics by
construction.